\section{Experimental Setup}
\label{sec:setup}

\subsection{Dataset}

We use the Human Activities and Postural Transitions (\hapt{})
dataset~\cite{reyes2015transition}, an extension of UCI
HAR~\cite{anguita2013public}. \hapt{} contains tri-axial linear
acceleration recorded at 50~Hz from a smartphone worn at the waist
by 30 subjects performing six basic activities:
\textsc{walking}, \textsc{upstairs}, \textsc{downstairs},
\textsc{sitting}, \textsc{standing}, and \textsc{laying}.
We adopt the canonical subject-disjoint train/validation/test split
of~\cite{reyes2015transition}, yielding 7{,}352 training windows,
1{,}515 validation windows, and 3{,}399 test windows of length~128
(2.56~s at 50~Hz). All accuracy and F1 numbers reported in this
paper are on the held-out test split.

For streaming evaluation we feed each test window sample-by-sample
at 50~Hz and reset the hidden state at every window boundary,
matching the deployed runtime exactly.

\subsection{Training Protocol}

Models are trained in PyTorch~2.x with the
Adam optimizer ($\eta{=}10^{-3}$, batch size~64, 100 epochs)
on a desktop CPU. The IHT sparsification mask follows the cubic
schedule of Section~\ref{sec:method:iht}, reaching the target
sparsity at epoch~50 and remaining frozen for the subsequent 50
epochs of fine-tuning. Each reported configuration is repeated
across the five seeds $\{0, 1, 2, 3, 4\}$; statistics are
mean~$\pm$~standard deviation unless otherwise noted.

\subsection{Deployment Hardware}

Two bare-metal targets are evaluated:
\begin{itemize}
\item \textbf{\arduino{}} (ATmega328P): 8-bit AVR core at
16~MHz; 32~KB Flash; 2~KB SRAM; hardware $8{\times}8$ multiplier;
\emph{no} floating-point unit. Toolchain: Arduino~IDE~2.3.x
with \texttt{avr-gcc} at \texttt{-Os}.
\item \textbf{\msp{}}: 16-bit MSP430 core at the calibrated
16~MHz DCO; 16~KB Flash; 512~B SRAM;
\emph{no hardware multiplier of any kind}; no FPU. Toolchain:
Code Composer Studio~12.x with the TI \texttt{msp430-elf-gcc}
bare-metal build, compiled at \texttt{-O2}. The Energia runtime
is \emph{not} used.
\end{itemize}

In both cases the deployed image consists of (i)~the
\texttt{fastgrnn.cpp} inference engine,
(ii)~the Q15 weight table and per-tensor scales
(\texttt{model\_weights.h}),
(iii)~the 256-entry sigmoid and tanh look-up tables, and
(iv)~a UART driver for streaming class labels at 9600 (\msp{}) or
115{,}200~baud (\arduino{}). All other Flash is empty.

\subsection{Cross-Platform Verification Protocol}
\label{sec:setup:verify}

To verify the claim of 100\% prediction agreement between the
FP32 PyTorch reference and the Q15 C implementation
(Section~\ref{sec:results:crossplatform}), we run both inference
pipelines on identical inputs and compare outputs. A Python harness
(\texttt{test\_inference\_python.py}) loads the deployed Q15
weights, executes a C-equivalent forward pass in NumPy (mirroring
the LUT activations and per-tensor dequantization steps), and
compares \texttt{argmax} predictions against the FP32 reference
for every test window. The C implementation on Arduino and MSP430
emits per-window predictions over UART, which are then compared
against the same reference. All three execution paths -- FP32
PyTorch, NumPy C-equivalent, and bare-metal C -- agree on all
3{,}399 windows.

\subsection{Live Sensor Configuration}

For live demonstrations a GY-521 module exposing the InvenSense
MPU-6050 is connected via I\textsuperscript{2}C. On the \msp{} we
drive USCI\_B0 directly (P1.6~$=$~SCL, P1.7~$=$~SDA, no Energia
\texttt{Wire} layer); on the \arduino{} the standard Wire library
is used. Only the accelerometer at $\pm 2$~g range is read,
matching the FP32 training-data scaling. A software timer paces
the sampling loop to 50~Hz, which leaves 7--11~ms of headroom per
sample after the inference step
(Section~\ref{sec:results:streaming}).
